% 第4章 4.1.1：(a) 光滑凸起 psi(x)=exp(-1/(1-x^2))（|x|<1），psi(0)=e^{-1}≈0.368，支撑[-1,1]；
% (b) 光滑截断 chi=1_{[-0.6,0.6]}*rho_{0.3}：在 K=[-0.3,0.3] 上为 1，支撑[-0.9,0.9]，含于 U=(-1,1)。
% 数据由同目录 generate.py 计算（bump.dat、cutoff.dat）。
\begin{tikzpicture}[font=\figurefont\footnotesize]
\pgfplotsset{every axis/.style={width=56mm,height=36mm,axis lines=left,
  x axis line style={draw=FigureStroke,line width=.6pt,-{Latex[length=1.6mm]}},
  y axis line style={draw=FigureStroke,line width=.6pt,-{Latex[length=1.6mm]}},
  tick style={draw=FigureStroke},tick label style={font=\figurefont\scriptsize,text=FigureInk},
  xlabel={$x$},xlabel style={at={(axis description cs:1,0)},anchor=west,font=\figurefont\scriptsize},
  title style={font=\figurefont\footnotesize,yshift=1mm},clip=false}}
\begin{axis}[name=A,xmin=-1.6,xmax=1.7,ymin=0,ymax=0.5,xtick={-1,0,1},ytick={0,0.368},
  yticklabels={$0$,$e^{-1}$},title={(a) 光滑凸起$\psi$}]
\addplot[draw=FigureBlue,line width=1.1pt] table[x=x,y=y] {figures/02-abstract-mathematics/03-function-spaces/tikz/c4-bump-cutoff/bump.dat};
\draw[FigureMuted,line width=.6pt,decorate,decoration={brace,mirror,amplitude=2.5pt}] (axis cs:-1,-0.13) -- node[below=2pt,font=\FigureNoteFont\scriptsize,text=FigureMuted]{$\operatorname{supp}\psi=[-1,1]$} (axis cs:1,-0.13);
\end{axis}
\begin{axis}[at={(A.east)},anchor=west,xshift=10mm,xmin=-1.3,xmax=1.4,ymin=0,ymax=1.25,
  xtick={-1,-0.3,0.3,1},xticklabels={$-1$,$-0.3$,$0.3$,$1$},ytick={0,1},title={(b) 光滑截断$\chi$}]
\fill[FigureYellowFill] (axis cs:-0.3,0) rectangle (axis cs:0.3,1);
\addplot[draw=FigureBlue,line width=1.1pt] table[x=x,y=y] {figures/02-abstract-mathematics/03-function-spaces/tikz/c4-bump-cutoff/cutoff.dat};
\node[font=\figurefont\scriptsize,text=FigureInk] at (axis cs:0,0.5) {$K$};
\draw[FigureGreen,line width=.8pt,{Parenthesis[]}-{Parenthesis[]}] (axis cs:-1,-0.33) -- node[below=1pt,font=\figurefont\scriptsize,text=FigureGreen]{$U$} (axis cs:1,-0.33);
\end{axis}
\end{tikzpicture}
